\nwuFrontHeading{摘\quad 要}

秦岭是我国重要的地理分界线和生态安全屏障，复杂的地形、气候和植被类型共同塑造了丰富的生物多样性。面向长期监测和精细化保护需求，来自样线调查、红外相机、遥感影像、物候记录和文献资料的数据往往具有尺度差异、缺测片段、分类口径不一致等问题。本文以公开场景构建博士论文示例，围绕秦岭生物多样性数据融合、关键变量组织、保护分区识别和管理解释四个环节展开，用于验证西北大学博士论文模板的排版能力。

本文首先梳理多源生态数据的来源、粒度和质量控制流程，将物种记录、栖息地因子、环境扰动和保护管理信息组织为可复核的数据层；其次建立跨尺度特征融合框架，处理样点观测与栅格环境变量之间的空间配准问题；再次构建保护优先区识别流程，综合物种丰富度、连通性、扰动压力和不确定性指标形成分区建议；最后给出结果解释与模板验证说明，展示章节标题、图表、公式、参考文献和博士论文后置材料在统一构建链路中的呈现效果。示例正文、数据和人名均为公开演示材料，不对应真实学位论文或保密项目。

排版层面，本模板依据《西北大学研究生学位论文规范》（研字〔2019〕7 号）和研究生院公开 Word 模板整理，覆盖中文封面、英文题名页、知识产权声明书、独创性声明、中英文摘要、目录、章节正文、结论、参考文献、附录、攻读博士学位期间取得的研究成果、致谢和作者简介。用户可在不修改公共样式的前提下替换 `extraTex/` 中的示例正文和 `meta.tex` 中的元信息。

\nwuKeywordsZhLine
